\section{Билет 14. Метод отсекающей гиперплоскости}

\ticketq{Метод отсекающей гиперплоскости.}

\paragraph{Постановка.}
\[
  \varphi_0(x) = c\T x \to \min , \qquad x \in \Omega ,
  \qquad \Omega = \{\, x \in \R^n \mid \varphi(x) \le 0 \,\} ,
\]
$\varphi$ выпукла, $\Omega \ne \varnothing$ — компакт. \textbf{Идея:}
заменить сложное $\Omega$ простым объемлющим многогранником и последовательно
добавлять линейные ограничения, отсекающие недопустимые точки.

\textbf{Приведение к этому виду.} Система $\varphi_i(x) \le 0$,
$i = \overline{1,m}$, сворачивается в одно ограничение
$\varphi(x) = \max_i \varphi_i(x)$ (максимум конечного числа выпуклых функций
выпуклен). Нелинейная выпуклая цель заменяется линейной вводом переменной:
$x^{(n+1)} \to \min$ при $\varphi_0(x) - x^{(n+1)} \le 0$.

\paragraph{Отсекающая гиперплоскость.}
Гиперплоскость $a\T x = \beta$ называется \textbf{отсекающей} для $y \notin \Omega$,
если
\[
  a\T x \le \beta \ \ \forall x \in \Omega
  \qquad\text{и}\qquad
  a\T y > \beta ,
\]
то есть новое ограничение сохраняет всё допустимое множество и выбрасывает $y$.

\textbf{Субградиент.} Вектор $p$ — субградиент выпуклой $\varphi$ в точке $y$,
если $\varphi(x) \ge \varphi(y) + p\T(x-y)$ $\forall x$. Для
дифференцируемой $\varphi$ субградиент совпадает с $\nabla\varphi(y)$.

\paragraph{Теорема (построение отсечения).}
\textit{Пусть $\varphi$ выпукла, $p \in \partial\varphi(y)$ и
$\varphi(y) > 0$. Тогда}
\[
  \varphi(y) + p\T(x - y) \le 0
\]
\textit{— отсекающее неравенство для точки $y$.}

\textsc{Док-во.} \emph{Допустимые точки сохраняются:} для $x \in \Omega$ по
определению субградиента
\[
  \varphi(y) + p\T(x-y) \ \le\ \varphi(x) \ \le\ 0 .
\]
\emph{Точка $y$ отсекается:} подстановка $x = y$ даёт
$\varphi(y) + p\T(y-y) = \varphi(y) > 0$, то есть неравенство нарушено.
$\blacksquare$

\paragraph{Внешние приближения.}
Строится последовательность многогранников
$S_0 \supset S_1 \supset \dots \supset \Omega$. На $k$-м шаге решается
линейная задача на $S_k$, её решение $x_k$; если $\varphi(x_k) > 0$, строится
субградиент $a_{k+1} \in \partial\varphi(x_k)$ и, при
$b^{(k+1)} = -\varphi(x_k) + a_{k+1}\T x_k$,
\[
  \varphi(x) \ \ge\ a_{k+1}\T x - b^{(k+1)} ,
  \qquad
  S_{k+1} = S_k \cap \big\{\, a_{k+1}\T x - b^{(k+1)} \le 0 \,\big\} .
\]

\textbf{Лемма (вложенность).} \textit{$\Omega \subset S_k$ и отсечение
корректно $\Rightarrow$ $\Omega \subset S_{k+1} \subset S_k$.}

\textsc{Док-во.} $x \in \Omega$ $\Rightarrow$ $x \in S_k$ (по условию) и
$a_{k+1}\T x - b^{(k+1)} \le \varphi(x) \le 0$ $\Rightarrow$ $x$ лежит в
пересечении, то есть в $S_{k+1}$. Текущая же $x_k$ выбрасывается:
$a_{k+1}\T x_k - b^{(k+1)} = \varphi(x_k) > 0$. $\blacksquare$

\paragraph{Алгоритм.}
\begin{csteps}
  \item Построить объемлющий многогранник-компакт $S_0 \supset \Omega$,
        заданный линейными ограничениями; $k = 0$.
  \item Решить $c\T x \to \min$, $x \in S_k$; получить $x_k$.
  \item $\varphi(x_k) \le 0$ — точка допустима, она и есть решение, конец.
  \item Иначе взять $a_{k+1} \in \partial\varphi(x_k)$, вычислить
        $b^{(k+1)} = -\varphi(x_k) + a_{k+1}\T x_k$ и положить
        $S_{k+1} = S_k \cap \{a_{k+1}\T x \le b^{(k+1)}\}$.
  \item $k := k+1$; останов при $\norm{x_{k+1} - x_k} < \eps$, иначе
        к~шагу~2.
\end{csteps}

\paragraph{Особенности.}
\textbf{Линейная подзадача.} От шага к шагу меняется только \emph{число}
ограничений $\Rightarrow$ применим симплекс-метод.

\textbf{Двойственность.} Добавление ограничения к прямой задаче — это
добавление одной \emph{переменной} к двойственной. Поэтому предыдущее
двойственное решение, дополненное нулевой компонентой, остаётся допустимым и
служит готовым начальным приближением: симплекс-метод продолжает работу, а не
стартует заново. В этом главный вычислительный выигрыш.

\paragraph{Теорема о сходимости.}
\textit{Пусть $\varphi$ непрерывна и выпукла, $\Omega$ и $S_0$ компактны,
субградиенты ограничены: $\norm{a_k} \le R$. Тогда любая предельная точка
$\{x_k\}$ — оптимальное решение, причём $\varphi(x_k) \to 0$.}

\textsc{Док-во.} $S_{k+1} \subseteq S_k \subseteq S_0$ и $S_0$ компактно
$\Rightarrow$ предельные точки существуют.

\emph{1)} Если при каком-то $k$ оказалось $\varphi(x_k) \le 0$, то
$x_k \in \Omega \subseteq S_k$ и $\varphi_0(x_k) = \min_{S_k}\varphi_0 \le
\varphi_0(x)$ $\forall x \in \Omega$ — точка уже оптимальна.

\emph{2)} Пусть $\varphi(x_k) > 0$ всегда и $x_k \to \bar{x}$ по подмножеству
$K$. Для $j < k$ имеем $x_k \in S_{j+1}$, то есть
$a_{j+1}\T(x_k - x_j) + \varphi(x_j) \le 0$, откуда
\[
  \varphi(x_j) \ \le\ \norm{a_{j+1}}\,\norm{x_j - x_k}
   \ \le\ R\,\norm{x_j - x_k} .
\]
Для $j,k \in K$ расстояние $\norm{x_j - x_k} \to 0$ $\Rightarrow$
$\varphi(x_j) \to 0$, и по непрерывности $\varphi(\bar{x}) = 0$, то есть
$\bar{x} \in \Omega$.

\emph{3)} $\varphi_0(x_k) = \min_{S_k}\varphi_0 \le \varphi_0(x)$
$\forall x \in \Omega$; переходя к пределу,
$\varphi_0(\bar{x}) \le \varphi_0(x)$ — точка $\bar{x}$ оптимальна.
$\blacksquare$

Содержательно: все $x_k$ недопустимы, поэтому дают \emph{нижние} оценки
оптимума, а предельная точка допустима и даёт верхнюю — метод зажимает оптимум
с двух сторон.

\paragraph{Ограничения.}
Может потребоваться много отсечений; подзадача растёт с каждой итерацией; при
слабых отсечениях область сокращается медленно; требуется выпуклость — для
невыпуклого $\Omega$ касательная гиперплоскость может отсечь часть допустимого
множества; при недифференцируемой $\varphi$ нужно уметь строить субградиенты.

\seefull{гл.~3, \S2 и \S7 (численный пример в \S7)}
